\subsection{One multiplicative-site cap for a primitive source}
\begin{statement}{P3F.40}{Proved}
Take any of the ten primitive-translation sources of P3K.22,
with its specified line, simple whole zero and ordinary remaining
source values. Replace one nonzero-linearization multiplicative
fibre by an actual transitive root-line cap. After the exact
finite-kernel cancellation of P3K.19, write its primitive norm
pair as $(m,n)$, with $m>0$, and retain its elliptic and delta
collars. The degree-two Mayer--Vietoris cokernel of its circle
orbit is exactly
\[
                         (\mathbb Z/m)^2.              \tag{F40.1}
\]
For $n\ne0$ this is $H_2(Y;\mathbb Z)$, and
$\pi_1(Y)=\pi_1(X)=\mathbb Z/|n|$. The sphere members among
these actual one-site constructions have exactly $m=1$ and
$n=\pm1$. The assertion concerns one multiplicative replacement;
it does not identify simultaneous replacements or their returns.
It holds for every allowed one-vertex cusp fan.
\end{statement}
\begin{proof}
\par\medskip\noindent\textbf{The marked boundary and its ordinary comparison.}
Let the fibre be $I_N$. Its marked generator has Tate valuation
$c_G$ prime to $N$. The physical outward suspension is the
inverse of the contracting lift and covers $+G$. Set $c=-c_G$.
In the actual boundary generators $a,b,\alpha,\gamma$, where
$a$ is the vanishing circle, $b$ the graph period and $\alpha$
the positive base meridian, the relations are
\[
 [\alpha,b]=a^N,\qquad[\alpha,\gamma]=a^{-c},\qquad
 [a,b]=[a,\gamma]=[b,\gamma]=1.                       \tag{F40.2}
\]
Indeed the Tate periods have first coordinates $Ns$ and
$c_Gs$ plus single-valued units, with $s=\log q/(2\pi i)$.
The indicated commutators are their actual changes under
$s\mapsto s+1$. Thus the boundary $E$ is the $a$-circle bundle
over $T^3_{b,\alpha,\gamma}$ with Euler form
$d\alpha\wedge(N\,db-c\,d\gamma)$.

Equivalently it is the three-torus mapping torus with matrix
\[
 J=\begin{pmatrix}1&N&-c\\0&1&0\\0&0&1\end{pmatrix}
\]
over $\alpha$. In its degree-two exterior-square quotient the
sole relation is
$c(a\wedge b)+N(a\wedge\gamma)=0$.
Choose integers $x,y$ with $Nx-cy=1$. The following are
actual integral degree-two generators:
\[
 \begin{split}
 z&=x(a\wedge b)+y(a\wedge\gamma),\qquad
 R=a\wedge\alpha,\\
 B&=b\wedge\gamma,\qquad
 L=(cb+N\gamma)\wedge\alpha.
 \end{split}                                                \tag{F40.3}
\]
The two suspended loops in (F40.3) are fixed by $J$, so their
tori are actual cycles. The change from the relation and $z$
to $(a\wedge b,a\wedge\gamma)$ has determinant $-1$.
Consequently $H_2(E)=\mathbb Z^4$, with
\[
 F=Nz,\qquad a\wedge\gamma=-cz.
\]
For boundary orientation $(a,b,\gamma,\alpha)$ its intersection
matrix in the displayed order is
\[
 Q_E=\begin{pmatrix}0&0&0&1\\0&0&1&0\\0&1&0&0\\1&0&0&0
 \end{pmatrix}.                                           \tag{F40.4}
\]
These signs follow by wedging the actual fibre and suspended
cycles: $zL=Nx-cy=1$ and $RB=1$.

Let $P$ be the complement of this fibre tube and
$W=T_{t_G}P$, using the physical outward suspension. The
component lattice is primitive. In fact it is a primitive
direct sublattice of the full trivial lattice, and the latter
has torsion-free quotient in $\operatorname{NS}(S)$ because
$\operatorname{MW}(S)=\mathbb Z$. Excision and the unimodular
source form give $H_1(P)=0$. The elliptic fibration with section
over the complementary disc presents its fundamental group as
a quotient of the abelian elliptic lattice; hence $\pi_1(P)=0$
and $\pi_1(W)=\mathbb Z\gamma$.

Denote the unchanged ordinary cap by $U$. Translation cycles
all $N$ components of its central fibre and preserves its
oriented graph circle. Its integral Wang sequence gives
$H_1(U)=\mathbb Z^2$, $H_2(U)=\mathbb Z^2$ and $H_3(U)=\mathbb Z$.
Write $C$ for the component-orbit class and $T$ for a lift of
graph times suspension. Then $F\mapsto NC$. The invariant sum
of the components maps to $F$ in the source Wang sequence, so
\[
 H_3(U)\longrightarrow H_3(T_{t_G}S)=\mathbb ZF
\]
is onto with degree one. The section generator of
$H_2(T_{t_G}S)=\mathbb Z[O]$ has boundary the primitive
meridian $\alpha$, since it meets the removed fibre once.
In the ordinary cap $\alpha$ bounds its section disc, while
$b,\gamma$ give the two primitive generators. These two
facts in the full-source Mayer--Vietoris sequence prove that
\[
 H_2(E)\longrightarrow H_2(W)\oplus H_2(U)              \tag{F40.5}
\]
is an integral isomorphism. In particular $H_2(W)=\mathbb Z^2$.
The ordinary cap map is $z\mapsto C$, $R,L\mapsto0$ and
$B\mapsto T$, after choosing the lift $T$. Thus the exterior
images of $R,L$ form an integral basis.

\par\medskip\noindent\textbf{The actual line frame and the zero fibre class.}
Write $M=\mathcal O(dG-O)$ as in P3K.22. Its degree on each
component of this retained fibre is zero. Its actual Tate Picard
representative $P(q)$ is a nonvanishing holomorphic unit, and
its line transition is
\[
                    (Z,V)\sim(q_TZ,P(q)V).
\]
Choose an analytic logarithm of $P(q)$. The expression
\[
 V\exp\left(-\log P(q)\frac{\log|Z|}{\log|q_T|}\right)
                                                               \tag{F40.6}
\]
is unchanged by this transition and gives a nowhere-zero smooth
frame of $M$ on the source boundary. A translation lift
$(Z,V)\mapsto(G(q)Z,\beta(q)V)$ acts in that frame by
\[
 \beta(q)\exp\left(-\log P(q)
                         \frac{\log|G(q)|}{\log|q_T|}\right).
                                                               \tag{F40.7}
\]
The degree-zero Tate charts have unit line transitions on every
component, so $\beta$ is the actual nonzero holomorphic unit
at this retained value. The exponential in (F40.7) has a
single-valued logarithm. Its graph and vanishing windings are
zero, and its base winding is zero as well.

Over the complementary disc, divide the specified degree-one
Hom section by an equation $g$ for its sole simple zero. This
gives an invertible lift on $P$. In the frame (F40.6) its ratio
differs from (F40.7) only by $g^{-1}$, whose boundary winding
is one signed base meridian. It has no graph winding. Form
the relative class $c_1(M,\text{frame})\in H^2(P,\partial P)$
and let $K\in H_2(P)$ be its Poincare--Lefschetz dual. The
relative Chern connecting map therefore gives
\[
                        (t_G)_*K-K=\pm F.              \tag{F40.8}
\]
Here the dual of $d\alpha$ on the source boundary is exactly
its fibre $F$; a graph winding would instead have given a rim
term, and (F40.7) proves that winding is zero. Equation (F40.8)
is a compact relative-cycle identity, not an assertion just in
the closed source. Hence $F=0$ in $H_2(W)$. Its torsion-freeness
and $F=Nz$ give $i(z)=0$ in the exterior.

Use $i(R),i(L)$ as the exterior basis and write
$i(B)=e\,i(R)+f\,i(L)$. The kernel of a boundary inclusion
into an oriented compact five-manifold is isotropic for the
boundary intersection form: bounding relative three-chains
give this immediately by the boundary intersection identity.
Pair $z$ with $B-eR-fL$ in (F40.4); this gives $f=0$.
The square of $B-eR$ is then $-2e$, so $e=0$. The complete
actual exterior map is thus
\[
                 i_2=\begin{pmatrix}0&1&0&0\\0&0&0&1\end{pmatrix}.
                                                               \tag{F40.9}
\]

\par\medskip\noindent\textbf{The marked cap and every degree-two column.}
Retain its primitive meridian
\[
                     \mu=m\alpha+jb+n\gamma.          \tag{F40.10}
\]
An additional vanishing-circle coefficient has zero class in
$H_1(E)$ because $\gcd(N,c)=1$ in (F40.2). The actual delta
phase stays in the line lift and has no orbit column.
Put $A=Nj-cn$. In the Tate root the affine deck has component
displacement $-A$: its norm differs from $B^n$ on the Tate
cover by the period $-j$, exactly as (F40.10) records. The
suspension component displacement is $mc$. Their component
orbits in $\mathbb Z/(Nm)$ therefore have number
\[
           \gcd(Nm,A,mc)=\gcd(m,A).
\]
The assumed actual transitive cap gives $\gcd(m,A)=1$.

Its central-root Wang sequence again gives
$H_2(\mathrm{cap})=\mathbb ZC\oplus\mathbb ZT$.
The generic fibre specializes to all $Nm$ components, so
$z\mapsto mC$. Choose integers $p,q$ with $mp+nq=-j$.
The actual degree-one cap map, retaining these graph-path
choices, is
\[
 \alpha\mapsto p\eta-nh,\qquad b\mapsto\eta,
                    \gamma\mapsto q\eta+mh.           \tag{F40.11}
\]
It is the quotient by (F40.10); the primitive graph and outward
suspension mark its two target generators. Its exterior square
computes projection to the $T$ summand. Thus $B\mapsto mT$
and $L\mapsto AT$, while $z,R$ have zero $T$ projection.
Write the still undetermined $C$ coefficients of $R,B,L$
as $r,s,t$. The isotropic-kernel identity for this actual cap
map, using (F40.4) and clearing denominators in its kernel
cycles, gives
\[
                         r=-A,\qquad mt=As.            \tag{F40.12}
\]
Equivalently multiply its matrix by $Q_E^{-1}$ and its transpose;
the off-diagonal entry is $m(r+A)$ and the first diagonal
entry is $2(mt-As)$. Both vanish. Since $m,A$ are coprime,
$s=mk$ and $t=Ak$ for an actual integer $k$. The determinant-one
target change $T_{\rm new}=T+kC$ retains this coefficient and
makes the cap matrix
\[
                 j_2=\begin{pmatrix}m&-A&0&0\\0&0&m&A\end{pmatrix}.
                                                               \tag{F40.13}
\]
No exterior inversion has been imposed on a marked cap.

The complete degree-two map is $(i_2,-j_2)^t$.
Eliminating its exterior basis columns $R,L$ leaves exactly
$\operatorname{diag}(m,m)$ on the kernel basis $z,B$.
This proves (F40.1), including both torsion summands and all
integer collar coefficients.

In the degree-one boundary order $(b,\alpha,\gamma)$ the
complete matrix is
\[
 D_1=\begin{pmatrix}0&0&1\\-1&-p&-q\\0&n&-m\end{pmatrix},
                         \qquad\det D_1=-n.            \tag{F40.14}
\]
For $n\ne0$ it is injective, so Mayer--Vietoris identifies
$H_2(Y)$ with (F40.1). The cap central mapping torus has
fundamental group $\mathbb Z^2$, since its source graph is
a circle preserved in orientation. In van Kampen the exterior
kills $b,\alpha$, hence $\eta=1$, $h^n=1$ and
$\gamma=h^m$. Coprimality of $m,n$ gives precisely
$\pi_1(Y)=\mathbb Z/|n|$. For $n=0$, primitivity already
forces $m=1$, and the same presentation is infinite cyclic.

A sphere therefore requires $m=1,n=\pm1$: (F40.1) injects
into orbit $H_2$ for every numerator, and the fundamental group
gives the other restriction. For these two numerators all
degree-one and degree-two groups vanish, and the closed oriented
five-manifold duality gives the integral sphere groups. The
same smooth classification and actual primitive Hopf-meridian
Gysin/fixed-tube maps as P3K.22 identify the orbit with $S^5$
and the total space with the oriented standard $S^6$.
The Hopf meridian also identifies $\pi_1(X)$ with $\pi_1(Y)$.
P3F.35's actual blowup/flop chain, equal to the identity outside
the cusp, preserves these marked caps and the sphere criterion
for every allowed fan.

The executable certificate
\path{code/part3_root/multiplicative_cap_boundary.py}
constructs both mapping-torus differentials, the integral cycle
bases, their intersection forms, all cap and exterior matrices,
and the retained target change. It checks every primitive-source
multiplicative marking and the exact symbolic determinant and
Smith form. The relative Chern framing above has an independent
audit recorded with the statement.
\end{proof}
